\documentclass[a4paper,10pt,french]{article}
\usepackage{../../CLT2425}
\geometry{a4paper,tmargin=1cm,bmargin=0.5cm,lmargin=1cm,rmargin=1cm}
%\renewcommand{\familydefault}{\sfdefault}
% \DeclareSymbolFont{numbers}{T1}{cmss}{m}{n}
%%  \SetSymbolFont{numbers}{}{T1}{phv}{bx}{n}
%  \DeclareMathSymbol{0}\mathalpha{numbers}{"30}
%  \DeclareMathSymbol{1}\mathalpha{numbers}{"31}
%  \DeclareMathSymbol{2}\mathalpha{numbers}{"32}
%  \DeclareMathSymbol{3}\mathalpha{numbers}{"33}
%  \DeclareMathSymbol{4}\mathalpha{numbers}{"34}
%  \DeclareMathSymbol{5}\mathalpha{numbers}{"35}
%  \DeclareMathSymbol{6}\mathalpha{numbers}{"36}
%  \DeclareMathSymbol{7}\mathalpha{numbers}{"37}
%  \DeclareMathSymbol{8}\mathalpha{numbers}{"38}
%  \DeclareMathSymbol{9}\mathalpha{numbers}{"39}

\renewcommand{\baselinestretch}{1.2}


\setlength{\columnsep}{30pt}
\begin{document}
\pagestyle{empty}

\titrebis{0.45}{Feuille 6 - Colinéarité et démonstration}

\textbf{{\large Exercice 1.}}\\
\I Soit ABC un triangle.\\
On définit les points M et N par les égalités vectorielles suivantes : $\Vj{\text{AM}}=\dfrac14\Vj{AB}$ et $\Vj{BN}=\dfrac34 \Vj{BC}$.\\

Complète la figure sur Géogébra. Code Capytale : \texttt{36f4-4415754}


On va chercher à montrer que les droites (MN) et (BC) sont parallèles.\\
Pour ce faire, nous allons montrer que les vecteurs $\Vj{MN}$ et $\Vj{BC}$ sont colinéaires, c'est-à-dire qu'on va chercher $k$ tel que $\Vj{MN}=k\Vj{BC}$.

\begin{multicols}{2}
\textbf{Tout d'abord, exprimons $\Vj{MB}$ en fonction de $\Vj{AB}$ : }
 \begin{eqnarray*}
\Vj{MB} & = & \Vj{M...} + \Vj{AB} \\
 &=&  \ \dfrac{\ ...\ }{\ ...\ } \ \Vj {AB} + \Vj {AB} \hspace*{1cm} \\ 
  &=&  \ \dfrac{\ ...\ }{\ ...\ } \ \Vj {AB} 
\end{eqnarray*}

\begin{center}
\textbf{Justifications}.
\end{center}
\phantom{.} \dotfill

\medskip

\phantom{.} \dotfill

\medskip

\phantom{.} \dotfill\\
\end{multicols}

\begin{multicols}{2}
\textbf{Ensuite, }
 \begin{eqnarray*}
\Vj{MN} & = & \Vj{MB} + \Vj{$\overset{\phantom{gtef}}{.....}$} \\
 &=& ... \ \Vj{AB} + \dfrac{...}{...}\ \Vj{$\overset{\phantom{gtef}}{.....}$}\\
  &=& \dfrac{...}{...} \left( \Vj{$\overset{\phantom{gtef}}{.....}$} + \Vj{$\overset{\phantom{gtef}}{.....}$} \right)\\
    &=& \ \dfrac{...}{...} \ \Vj{$\overset{\phantom{gtef}}{.....}$}
\end{eqnarray*}


\begin{center}
\textbf{Justifications}.
\end{center}
\phantom{.} \dotfill

\medskip

\phantom{.} \dotfill

\bigskip

\phantom{.} \dotfill


\bigskip

\phantom{.} \dotfill
\\
\end{multicols}

\textbf{Conclusion} : les vecteurs ..... et .... sont colinéaires. 

\smallskip

\fbox{\textbf{Ainsi} : les droites ..... et ..... sont parallèles.}

\bigskip

\bigskip

\textbf{{\large Exercice 2.}}\\
\I Soit ABC un triangle.\\
On définit les points M et N par les égalités vectorielles suivantes : $\Vj{BM}=2\Vj{AC}$ et $\Vj{AN}=\dfrac{1}{2} \Vj{BC}$.\\
Réaliser une figure dans Géogébra via Capytale.

\medskip

\textbf{Observations} : \dotfill

\begin{multicols}{2}
\textbf{D'une part, }
\begin{eqnarray*}
\Vj{CM} & = & \Vj{...\ B} + \Vj{$\overset{\phantom{gtef}}{.....}$} \hspace*{1cm} (relation \ de \ ................. )\\
 &=& \Vj {CB} + 2\Vj {AC}
\end{eqnarray*}
\ \\
\ \\
\textbf{D'autre part, }
\begin{eqnarray*}
\Vj{CN} &=&  \Vj{...\ A} + \Vj{$\overset{\phantom{gtef}}{.....}$} \hspace*{1cm} (relation \ de \ ................. )\\
&=& -\Vj{$\overset{\phantom{gtef}}{.....}$} + \dfrac{1}{2}\Vj{$\overset{\phantom{gtef}}{.....}$} \\
&=& -\dfrac{1}{2}\Vj{$\overset{\phantom{gtef}}{.....}$}  - \Vj{$\overset{\phantom{gtef}}{.....}$}
\end{eqnarray*}
\end{multicols}

\bigskip


\I \begin{spacing}{1.8}
\I On remarque que : \fbox{$\Vj{CM} = \phantom{\dfrac{1}{2}} \ ....\  \Vj{CN}$} donc les vecteurs ...... et ....... sont ...............................
\medskip

\fbox{\textbf{Ainsi} les points .... , ..... et ..... sont ....................... .} 
\end{spacing}


\end{document}